\subsection{Removing every nontrivial singleton-pendant corner}

\begin{theorem}[Proved: arbitrary pendant-corner retraction]
\label{CP2:A-singleton-pendant-any-corner}
Let $p$ be a source or sink singleton pendant in a finite category
$\mathcal C$, in the whole-category sense of the preceding section.
Let $q$ be its neighbor, $L=\mathcal C(p,p)$, and
$N=\mathcal C(q,q)$. Replace $L$ by its identity, keeping all
other arrows and remaining composition values, to obtain
$\mathcal C'$. There are actual functors
\[
 i:\mathcal C'\hookrightarrow\mathcal C,
 \qquad R:\mathcal C\longrightarrow\mathcal C',
 \qquad Ri=\operatorname{id}_{\mathcal C'}.
\]
Every positive-degree product in the integral relative ring
$H^*(B\mathcal C,B\mathcal C';\mathbb Z)$ is zero.
Consequently, if $B\mathcal C\simeq\CP^2$, the inclusion
$Bi$ is a homotopy equivalence. The smaller category has
$|\operatorname{Mor}\mathcal C|-(|L|-1)$ arrows.
No hypothesis on the homology of $BL$ is needed.
\end{theorem}
\begin{proof}
It suffices to treat a source pendant. Write its sole cross
arrow as $f:p\to q$. There is no other strict hom-set incident
to $p$, in particular no arrow from $q$ to $p$.
Every right $L$ and left $N$ action on $\{f\}$ fixes $f$.
Sending all elements of $L$ to $1_p$ and fixing every other
arrow therefore preserves all products, including $f\ell=f$
and $nf=f$. The remaining arrows form a subcategory and give
the displayed inclusion and retraction.

Let $\mathcal E=\mathcal C\setminus\{p\}$ be the full deletion.
A simplex containing $p$ can contain no object other than
$p,q$, since any further strict object change would compose
to another strict arrow incident to $p$. Thus there are exact
unions of nerve realizations along CW subcomplexes
\[
 B\mathcal C=B\mathcal E\cup_{BN}J,
 \qquad B\mathcal C'=B\mathcal E\cup_{BN}K,
 \qquad J=BL*BN,
 \qquad K=\{\ell_0\}*BN.
\]
Here $\ell_0$ is the unique vertex of $BL$. The identifications
with joins are natural: a mixed simplex is precisely a string
of $L$ arrows, the unique cross arrow, and a string of $N$
arrows. The identity-corner inclusion is the vertex inclusion
$\{\ell_0\}\subset BL$, so $K$ is the actual cone subcomplex
of $J$ corresponding to $\mathcal C'$.

In particular $BN\subset K$, and collapsing the retained
subcategory gives
\[
 B\mathcal C/B\mathcal C'\cong J/K.
\]
The subcomplex $K$ is nonempty and contractible. A contraction
of $K$ extends across $J$ by the homotopy extension property,
showing that the quotient map $J\to J/K$ is a homotopy
equivalence. Since $BL$ and $BN$ are connected based CW
spaces, their join is homotopy equivalent to the suspension
of their smash product. Consequently
\[
 B\mathcal C/B\mathcal C'\simeq
 BL*BN\simeq\Sigma(BL\wedge BN).
\]
The reduced integral cohomology of a suspension has zero
products in positive degrees. The relative ring of the CW
pair is its quotient's reduced cohomology ring, proving the
relative-product assertion. Under $B\mathcal C\simeq\CP^2$,
Lemma~\ref{CP2:A-retract-relative-products} applies to the
actual inclusion and retraction and makes $Bi$ a homotopy
equivalence. Exactly the $|L|-1$ nonidentity corner arrows
were removed. Taking opposite categories proves the sink
case, without changing the space or the arrow count.
\end{proof}

\begin{corollary}[Proved: pendant corners at eighteen arrows]
\label{CP2:A-singleton-pendant-any-eighteen}
An eighteen-arrow category of type $\CP^2$ has a trivial
corner at every source or sink singleton pendant.
More generally, every singleton-pendant corner is trivial
in an arrow-minimal finite category of type $\CP^2$.
\end{corollary}
\begin{proof}
A nontrivial corner could be replaced by its identity by the
theorem, producing an actual model with fewer arrows. In the
eighteen-arrow case the result would have at most seventeen
arrows, contradicting Theorem~\ref{CP2:A-global-lower-eighteen}.
For an arrow-minimal model the reduction directly contradicts
minimality.
\end{proof}
